\chapter{1977年广东卷}

\section{解答题}

\begin{problem}
计算.
\begin{enumerate}
    \item
    \(\left(-3\right)^2-\left(-1\frac{1}{2}\right)^3\times\frac{2}{9}-6\div\left|-\frac{2}{3}\right|\).
    \item
    \(\frac{x^2+2x+1}{x^3-x}\cdot\frac{x}{x+1}-\frac{1}{x+1}\).
    \item
    \(\frac{\sqrt{20}}{8}-\sqrt{\frac{9}{5}}+\frac{5}{\sqrt{45}}\).
\end{enumerate}
\end{problem}

\begin{solution}
\begin{enumerate}
\item 因为 \(\left(-3\right)^2=9\)，\(\left(-1\frac{1}{2}\right)^3=\left(-\frac{3}{2}\right)^3=-\frac{27}{8}\)，且 \(\left|-\frac{2}{3}\right|=\frac{2}{3}\)，所以原式
\[
9-\left(-\frac{27}{8}\right)\times\frac{2}{9}-6\div\frac{2}{3}=9+\frac{3}{4}-9=\frac{3}{4}
\]

\item 原式有意义要求 \(x^3-x\ne0\) 且 \(x+1\ne0\). 因为 \(x^3-x=x\left(x-1\right)\left(x+1\right)\)，所以定义条件为 \(x\ne-1\)、\(x\ne0\) 且 \(x\ne1\). 在此条件下，原式
\[
\begin{gathered}
\frac{\left(x+1\right)^2}{x\left(x-1\right)\left(x+1\right)}\cdot\frac{x}{x+1}-\frac{1}{x+1} \\
=\frac{1}{x-1}-\frac{1}{x+1}=\frac{\left(x+1\right)-\left(x-1\right)}{\left(x-1\right)\left(x+1\right)}=\frac{2}{x^2-1}.
\end{gathered}
\]
因此化简结果为 \(\frac{2}{x^2-1}\)，定义条件仍为 \(x\ne-1\)、\(x\ne0\) 且 \(x\ne1\).

\item 因为 \(\sqrt{20}=2\sqrt{5}\)，所以 \(\frac{\sqrt{20}}{8}=\frac{\sqrt{5}}{4}\)；又 \(\sqrt{\frac{9}{5}}=\frac{3\sqrt{5}}{5}\)，且
\(\frac{5}{\sqrt{45}}=\frac{5}{3\sqrt{5}}=\frac{\sqrt{5}}{3}\). 因此原式
\[
\frac{\sqrt{5}}{4}-\frac{3\sqrt{5}}{5}+\frac{\sqrt{5}}{3}
=\frac{15\sqrt{5}-36\sqrt{5}+20\sqrt{5}}{60}=-\frac{\sqrt{5}}{60}
\]
\end{enumerate}
\end{solution}

\begin{problem}
解方程：\(\frac{x+3}{x+2}+\frac{x+1}{x-2}-1=0\)，并检验所得的结果是否是原方程的根.
\end{problem}

\begin{solution}
原方程的分母不能为零，所以必须有 \(x\ne-2\) 且 \(x\ne2\). 在这个条件下，\(\left(x+2\right)\left(x-2\right)\ne0\)，原方程两边同乘这个非零因式，得
\[
\left(x+3\right)\left(x-2\right)+\left(x+1\right)\left(x+2\right)-\left(x+2\right)\left(x-2\right)=0
\]
展开并合并同类项，得
\[
x^2+x-6+x^2+3x+2-x^2+4=x^2+4x=x\left(x+4\right)=0
\]
所以候选值为 \(x=0\) 或 \(x=-4\). 两者都不等于限制值 \(-2\) 和 \(2\)，代入原方程时各分母均非零；同时，上述同乘非零因式的变形是等价变形. 因此 \(x=0\) 和 \(x=-4\) 都是原方程的根，原方程的解为 \(x=0\)，\(x=-4\).
\end{solution}

\begin{problem}
计算.
\begin{enumerate}
    \item
    \(\left(\frac{1}{2}\right)^{-1}-4\times\left(-2\right)^{-3}+\left(\frac{1}{2}\right)^0-9^{-\frac{1}{2}}\).
    \item
    \(2\lg\left(a^{\frac{1}{3}}b^2\right)-\lg\frac{\sqrt{b}}{8}-3\lg2\)，其中 \(\lg a=0.33\)，\(\lg b=0.22\).
\end{enumerate}
\end{problem}

\begin{solution}
\begin{enumerate}
\item 由负整数指数幂和零指数幂的定义，\(\left(\frac{1}{2}\right)^{-1}=2\)，\(\left(-2\right)^{-3}=-\frac{1}{8}\)，\(\left(\frac{1}{2}\right)^0=1\)，\(9^{-\frac{1}{2}}=\frac{1}{\sqrt{9}}=\frac{1}{3}\). 所以原式
\[
2-4\times\left(-\frac{1}{8}\right)+1-\frac{1}{3}=2+\frac{1}{2}+1-\frac{1}{3}=\frac{19}{6}=3\frac{1}{6}
\]

\item 对数式有意义时，必须有 \(a>0\) 且 \(b>0\). 由对数的积、商和幂运算法则，\(\lg8=3\lg2\)，从而原式
\[
\begin{gathered}
2\left(\frac{1}{3}\lg a+2\lg b\right)-\left(\frac{1}{2}\lg b-\lg8\right)-3\lg2 \\
=\frac{2}{3}\lg a+4\lg b-\frac{1}{2}\lg b+3\lg2-3\lg2 \\
=\frac{2}{3}\lg a+\frac{7}{2}\lg b.
\end{gathered}
\]
代入 \(\lg a=0.33\) 和 \(\lg b=0.22\)，得原式
\[
\frac{2}{3}\times0.33+\frac{7}{2}\times0.22=0.22+0.77=0.99
\]
\end{enumerate}
\end{solution}

\begin{problem}
证明恒等式：\(\left(\frac{1}{\cos x}+1\right)\left(\frac{1}{\cos x}-1\right)\cot x=\tan x\).
\end{problem}

\begin{solution}
左边有意义要求 \(\cos x\ne0\) 且 \(\sin x\ne0\)，右边的 \(\tan x\) 也要求 \(\cos x\ne0\). 因此在等式两边都有意义的范围内，可以使用 \(\cot x=\frac{\cos x}{\sin x}\) 和 \(1-\cos^2x=\sin^2x\). 于是左边
\[
\begin{gathered}
\left(\frac{1}{\cos x}+1\right)\left(\frac{1}{\cos x}-1\right)\cot x
=\frac{1+\cos x}{\cos x}\cdot\frac{1-\cos x}{\cos x}\cdot\frac{\cos x}{\sin x} \\
=\frac{1-\cos^2x}{\cos x\sin x}=\frac{\sin^2x}{\cos x\sin x}=\frac{\sin x}{\cos x}=\tan x.
\end{gathered}
\]
所以在原定义域内左边恒等于右边，恒等式得证.
\end{solution}

\begin{problem}
已知小山上的电视发射塔 \(AB\) 的高是 \(50\text{ 米}\)，自山下地面的 \(C\) 点测得塔底 \(B\) 的仰角是 \(40^\circ\)，塔顶 \(A\) 的仰角是 \(70^\circ\)，求小山 \(BD\) 的高. （精确到 \(0.01\text{ 米}\)，且 \(\sin40^\circ=0.643\)，\(\sin20^\circ=0.342\)）

\bitmapfigure[max width=6.0cm]{img/1977/guangdong_hainan_included/q5_fig1.png}
\end{problem}

\begin{solution}
由图可知 \(A\)、\(B\)、\(D\) 共线，且 \(CD\) 为水平线，\(AD\perp CD\). 在直角三角形 \(ACD\) 中，\(\angle ACD=70^\circ\)，所以
\(\angle CAD=180^\circ-90^\circ-70^\circ=20^\circ\). 因为 \(A\)、\(B\)、\(D\) 共线，所以 \(\angle CAB=20^\circ\). 两条视线 \(CA\) 和 \(CB\) 与水平线的仰角分别为 \(70^\circ\) 和 \(40^\circ\)，故 \(\angle ACB=70^\circ-40^\circ=30^\circ\).

在三角形 \(ABC\) 中，由正弦定理，\(AB\) 对应 \(\angle ACB\)，\(BC\) 对应 \(\angle CAB\)，于是
\[
BC=\frac{AB\sin20^\circ}{\sin30^\circ}=\frac{50\sin20^\circ}{\frac{1}{2}}=100\sin20^\circ
\]
在直角三角形 \(BDC\) 中，\(\angle BCD=40^\circ\)，所以
\[
BD=BC\sin40^\circ=100\sin20^\circ\sin40^\circ=100\times0.342\times0.643=21.9906\text{ 米}
\]
按题意精确到 \(0.01\text{ 米}\)，得 \(BD\approx21.99\text{ 米}\). 因此小山 \(BD\) 的高约为 \(21.99\text{ 米}\).
\end{solution}

\begin{problem}
如图，已知 \(AE\) 和 \(AF\) 是 \(\odot O\) 的切线，\(M\)、\(N\) 是切点，\(EF\) 过圆心且垂直于 \(AO\). 求证：\(AE=AF\).

\bitmapfigure[max width=4.8cm]{img/1977/guangdong_hainan_included/q6_fig1.png}
\end{problem}

\begin{solution}
连接 \(OM\)、\(ON\). 因为 \(M\)、\(N\) 是切点，所以半径分别垂直于切线，即 \(OM\perp AE\)、\(ON\perp AF\). 因此三角形 \(OMA\) 和 \(ONA\) 都是直角三角形，且斜边 \(OA\) 是公共边，直角边 \(OM=ON\)（同圆半径相等），所以
\[
\triangle OMA\cong\triangle ONA
\]
从而 \(\angle MAO=\angle OAN\). 又因为 \(A\)、\(M\)、\(E\) 共线，\(A\)、\(N\)、\(F\) 共线，所以 \(\angle EAO=\angle OAF\).

由于 \(EF\perp AO\) 且 \(O\) 在 \(EF\) 上，\(\angle AOE=\angle AOF=90^\circ\). 在三角形 \(AOE\) 和 \(AOF\) 中，\(\angle EAO=\angle OAF\)，\(\angle AOE=\angle AOF\)，且 \(AO\) 为公共边，所以
\[
\triangle AOE\cong\triangle AOF
\]
因此对应边相等，得 \(AE=AF\).
\end{solution}

\begin{problem}
如果平面内一动点 \(P\) 到点 \(\left(-4,0\right)\) 的距离等于它到直线 \(x=-\frac{25}{4}\) 的距离的 \(\frac{4}{5}\)，
\begin{enumerate}
\item 求 \(P\) 点的轨迹方程；
\item 这又是一条什么曲线？
\end{enumerate}
\end{problem}

\begin{solution}
\begin{enumerate}
\item
设动点 \(P\) 的坐标为 \(\left(x,y\right)\). 根据两点间距离公式，\(P\) 到定点 \(\left(-4,0\right)\) 的距离为 \(\sqrt{\left(x+4\right)^2+y^2}\)；根据点到直线的距离公式，\(P\) 到直线 \(x=-\frac{25}{4}\) 的距离为 \(\left|x+\frac{25}{4}\right|\). 因此由题意得
\[
\sqrt{\left(x+4\right)^2+y^2}=\frac{4}{5}\left|x+\frac{25}{4}\right|
\]
等式两边都非负，所以将两边平方是等价变形，不会引入增根. 平方并整理，得
\[
\begin{gathered}
25\left(\left(x+4\right)^2+y^2\right)=16\left(x+\frac{25}{4}\right)^2,\\
25x^2+200x+400+25y^2=16x^2+200x+625,\\
9x^2+25y^2=225.
\end{gathered}
\]
两边除以 \(225\)，得到轨迹方程
\[
\frac{x^2}{25}+\frac{y^2}{9}=1
\]
\item 这个方程是椭圆的标准方程，且分母较大的 \(25\) 对应 \(x\) 项，所以它是以原点为中心、长半轴为 \(5\)、短半轴为 \(3\) 的椭圆.
\end{enumerate}
\end{solution}

\section{参考题}

\begin{problem}
方程组 \(\begin{cases}
x_1+3x_2+2x_3-x_4=1,\\
x_2-2x_3+x_4=-1,\\
2x_1+x_3+2x_4=0,\\
2x_2+ax_3-2x_4=1
\end{cases}\)
中，
\begin{enumerate}
\item 当 \(a\) 取什么值时，该方程组有解？（有多少组解？）
\item 又当 \(a\) 取什么值时，方程组无解？
\end{enumerate}
\end{problem}

\begin{solution}
将原方程组的四个方程依次记为前四个方程. 第一个方程乘 \(2\) 后减去第三个方程，得
\[
6x_2+3x_3-4x_4=2
\]
第二个方程乘 \(6\) 后减去这个等式，得
\[
-15x_3+10x_4=-8
\]
第二个方程乘 \(2\) 后减去第四个方程，得
\[
-\left(a+4\right)x_3+4x_4=-3
\]
将前一个等式乘 \(2\) 后减去后一个等式乘 \(5\)，消去 \(x_4\)，得到
\[
\left(5a-10\right)x_3=-1
\]

\begin{enumerate}
\item 当 \(a\ne2\) 时，\(5a-10=5\left(a-2\right)\ne0\)，所以
\[
x_3=-\frac{1}{5\left(a-2\right)}
\]
此时由 \(-15x_3+10x_4=-8\) 唯一确定 \(x_4\)，再由第二个方程唯一确定 \(x_2\)，最后由第三个方程唯一确定 \(x_1\). 令 \(d=a-2\)，可具体求得
\(x_4=\frac{13-8a}{10d}\)，\(x_2=\frac{3-2a}{10d}\)，\(x_1=\frac{4a-6}{5d}\).
因此当 \(a\ne2\) 时，方程组有且只有一组解.

\item 当 \(a=2\) 时，上面的两个关于 \(x_3\)、\(x_4\) 的等式变为
\[
\begin{cases}
-15x_3+10x_4=-8,\\
-6x_3+4x_4=-3.
\end{cases}
\]
第二式乘 \(\frac{5}{2}\) 后，左边变为 \(-15x_3+10x_4\)，而右边变为 \(-\frac{15}{2}\)，这与第一式右边 \(-8\) 不相等，产生矛盾. 所以当 \(a=2\) 时，方程组无解.
\end{enumerate}
\end{solution}

\begin{problem}
有一份印刷品，其排版面积（矩形）占 \(432\text{ 平方厘米}\)，它的两边都留有 \(4\text{ 厘米}\) 宽的空白，顶部和底部留有 \(3\text{ 厘米}\) 宽的空白. 问应该如何选取纸张的尺寸，才能使纸的用量最少？
\end{problem}

\begin{solution}
设纸张左右方向的长度为 \(x\text{ 厘米}\)，上下方向的长度为 \(y\text{ 厘米}\)，纸张面积为 \(S\text{ 平方厘米}\). 由于左右各留出 \(4\text{ 厘米}\)、上下各留出 \(3\text{ 厘米}\)，排版部分的长和宽分别为 \(x-8\) 和 \(y-6\). 因而 \(x>8\)、\(y>6\)，并且
\[
(x-8)(y-6)=432
\]
要使纸张面积 \(S=xy\) 最小，令 \(t=x-8\)，则 \(t>0\)，且
\(y=\frac{432}{t}+6\)，\(x=t+8\).
所以
\[
S=(t+8)\left(\frac{432}{t}+6\right)=\frac{3456}{t}+6t+480
\]
因为 \(t>0\)，所以 \(\frac{3456}{t}>0\) 且 \(6t>0\). 由基本不等式，
\[
S\ge2\sqrt{\frac{3456}{t}\cdot6t}+480=2\sqrt{20736}+480=768
\]
等号成立当且仅当
\[
\frac{3456}{t}=6t
\]
于是 \(t^2=576\). 因为 \(t>0\)，所以 \(t=24\)，从而
\(x=t+8=32\text{ 厘米}\)，\(y=\frac{432}{24}+6=24\text{ 厘米}\).
此时 \(S=32\times24=768\text{ 平方厘米}\)，达到由基本不等式得到的下界，故确为最小值. 所以纸张左右长取 \(32\text{ 厘米}\)、上下宽取 \(24\text{ 厘米}\) 时，用纸量最少，最小面积为 \(768\text{ 平方厘米}\).
\end{solution}
